% Part: computability
% Chapter: recursive-functions
% Section: bounded-minimization

\documentclass[../../../include/open-logic-section]{subfiles}

\begin{document}

\olfileid{cmp}{rec}{bmi}
\olsection{有界极小化}

\begin{explain}
将函数定义为满足某个性质或关系~$P$ 的最小数，这通常很有用。
如果 $P$ 是可判定的，我们可以简单地尝试所有可能的数 $0, 1, 2, \dots$ 来计算这个函数，直到找到满足~$P$ 的最小数。
这类\emph{无界搜索}使我们超出了原始递归函数的范畴。
然而，如果我们只关心\emph{小于某个独立给定界}的最小数，我们仍停留在原始递归的范畴内。
更一般地，假设我们有一个原始递归关系~$R(x,z)$。
考虑将 $x$ 和~$y$ 映射为满足 $R(x, z)$ 的最小 $z < y$ 的函数。
这个函数也可以通过测试 $R(x, 0)$、$R(x, 1)$、……、$R(x, y-1)$ 来计算。
但它为什么是原始递归的呢？
\end{explain}

\begin{prop}
如果 $R(\vec x, z)$ 是原始递归的，那么函数
$m_R(\vec{x}, y)$也是原始递归的，该函数在存在满足 $R(\vec x, z)$ 的 $z<y$ 时返回最小的 $z$，否则返回 $y$。
我们将函数~$m_R$ 记作
\[
\bmin{z < y}{R(\vec{x}, z)},
\]
\end{prop}

\begin{proof}
注意，不存在~$z < 0$ 使得 $R(\vec{x}, z)$ 成立，因为根本不存在 $z < 0$。
所以 $m_R(\vec x, 0) = 0$。

当上界形如 $y + 1$ 时，我们有三种情况：
\begin{enumerate}
\item 存在 $z < y$ 使得 $R(\vec{x}, z)$ 成立，此时
$m_R(\vec{x}, y+1) = m_R(\vec{x}, y)$。
\item 不存在这样的~$z<y$ 但 $R(\vec{x}, y)$ 成立，此时
$m_R(\vec{x}, y+1) = y$。
\item 不存在 $z < y+1$ 使得 $R(\vec{x}, z)$ 成立，此时
$m_R(\vec{z}, y+1) = y+1$。
\end{enumerate}
因此，我们可以如下通过原始递归来定义 $m_R(\vec x, y+1)$：
\begin{align*}
m_R(\vec{x}, 0) & = 0\\
m_R(\vec{x}, y+1) & =
\begin{cases}
m_R(\vec{x}, y) & \text{若 $m_R(\vec x, y) \neq y$}\\
y & \text{若 $m_R(\vec x, y) = y$ 且 $R(\vec{x}, y)$}\\
y+1 & \text{否则。}
\end{cases}
\end{align*}
注意，存在 $z<y$ 使得 $R(\vec x, z)$ 成立，当且仅当 $m_R(\vec x,
y) \neq y$。
\end{proof}

\begin{prob}
假设 $R(\vec x, z)$ 是原始递归的。
请从~$\Char{R}$ 出发通过原始递归定义函数
$m'_R(\vec{x}, y)$，该函数在存在满足 $R(\vec{x}, z)$ 的 $z<y$ 时返回最小的 $z$，否则返回 $0$。
\end{prob}

\end{document}